\documentclass[10pt,letterpaper]{article}
\usepackage{cornell-notes}

\title{Chapter 9: Security}
\author{}
\date{}

\setDocTitle{Chapter 9: Security}
\setDocSubtitle{Operating Systems}
\setDocVersion{v1.0}
\setDocStatus{Study Companion}
\setDocOwner{Operating Systems Cornell Notes}
\setDocAuthor{}
\setDocDate{\today}
\setDocSummary{Cornell-format study and review notes for Chapter 9: Security.}

\setCornellCollection{Operating Systems Cornell Notes}
\setCornellUnitType{Chapter}
\setCornellUnitNumber{9}
\setCornellUnitTitle{Security}
\setCornellCompanionLabel{Study and review companion}

\hypersetup{pdftitle={Chapter 9: Security},pdfsubject={Cornell notes},pdfauthor={}}

\begin{document}
\maketitle
\section*{Learning Objectives}
\begin{studybox}{By the end of this chapter, you should be able to}
\begin{enumerate}
  \item Classify security goals, threats, attackers, and operating-system trust assumptions.
  \item Model authorization with protection domains, access-control lists, and capabilities.
  \item Explain multilevel information-flow rules and covert channels.
  \item Distinguish secret-key encryption, public-key encryption, hashes, signatures, certificates, and trusted hardware.
  \item Compare knowledge-, possession-, and biometric-based authentication.
  \item Trace the causes and effects of common memory-safety, injection, and race-condition exploits.
  \item Distinguish insider attacks, Trojan horses, viruses, worms, spyware, and rootkits.
  \item Evaluate firewalls, malware scanning, code signing, jailing, intrusion detection, and mobile-code confinement as defense in depth.
\end{enumerate}
\end{studybox}

\begin{studybox}{Section Map}
\hyperlink{sec-9-1}{9.1} \quad \hyperlink{sec-9-2}{9.2} \quad \hyperlink{sec-9-3}{9.3} \quad \hyperlink{sec-9-4}{9.4} \quad \hyperlink{sec-9-5}{9.5} \quad \hyperlink{sec-9-6}{9.6} \quad \hyperlink{sec-9-7}{9.7} \quad \hyperlink{sec-9-8}{9.8} \quad \hyperlink{sec-9-9}{9.9} \quad \hyperlink{sec-9-10}{9.10} \quad \hyperlink{sec-9-11}{9.11} \quad \hyperlink{sec-9-12}{9.12}
\end{studybox}

\section*{Cornell Cue-and-Notes Area}
\small
\begin{longtable}{C N}
\rowcolor{navy}
\color{white}\textbf{Cue / Recall Prompt} &
\color{white}\textbf{Notes}\\
\endfirsthead
\rowcolor{navy}
\color{white}\textbf{Cue / Recall Prompt} &
\color{white}\textbf{Notes (continued)}\\
\endhead
\rowcolor{ice}\multicolumn{2}{l}{\hypertarget{sec-9-1}{}\textbf{Section 9.1: The Security Environment}}\\*\hline
\textbf{Security goals} & Confidentiality limits disclosure, integrity prevents unauthorized change, and availability keeps services usable. Authenticity, accountability, and privacy refine these goals.\\\hline
\textbf{9.1.1 Threats} & Threats include unauthorized information exposure, modification, fabrication, service interruption, and resource misuse. A policy states permitted behavior; a mechanism attempts to enforce it.\\\hline
\textbf{9.1.2 Attackers} & Attackers differ in access, skill, resources, persistence, and motive. Insider knowledge and legitimate credentials can bypass assumptions made for external threats.\\\hline
\rowcolor{ice}\multicolumn{2}{l}{\hypertarget{sec-9-2}{}\textbf{Section 9.2: Operating Systems Security}}\\*\hline
\textbf{9.2.1 Secure-system challenge} & A system must enforce policy across every path, resist implementation defects, and remain manageable. Complexity and compatibility enlarge the attack surface.\\\hline
\textbf{9.2.2 Trusted computing base} & The TCB is the hardware and software whose correct operation is required for security. A small reference monitor should mediate every access, resist tampering, and be analyzable.\\\hline
\rowcolor{ice}\multicolumn{2}{l}{\hypertarget{sec-9-3}{}\textbf{Section 9.3: Controlling Access to Resources}}\\*\hline
\textbf{Access matrix} & Rows represent protection domains, columns represent objects, and cells contain allowed operations. Real systems store sparse row- or column-oriented representations.\\\hline
\textbf{9.3.1 Protection domains} & A domain is a set of access rights associated with a user, process, procedure, or execution mode. Controlled domain switching changes authority.\\\hline
\textbf{9.3.2 ACLs} & An access-control list attaches subjects and allowed operations to each object. It answers who may use this object and makes object-wide review or revocation direct.\\\hline
\textbf{9.3.3 Capabilities} & A capability is an unforgeable reference combining an object identity with rights. Capability lists answer what this domain can use; delegation is natural, but global revocation needs additional mechanisms.\\\hline
\rowcolor{ice}\multicolumn{2}{l}{\hypertarget{sec-9-4}{}\textbf{Section 9.4: Formal Models of Secure Systems}}\\*\hline
\textbf{9.4.1 Multilevel security} & Information-flow models assign classifications and clearances. Confidentiality rules prevent reading above clearance and writing information to a lower level; integrity models reverse the emphasis to protect trustworthy data.\\\hline
\textbf{9.4.2 Covert channels} & A covert channel encodes information through a resource not intended for communication, such as timing, CPU load, file locks, or storage allocation. Closing every shared modulation path is difficult.\\\hline
\rowcolor{ice}\multicolumn{2}{l}{\hypertarget{sec-9-5}{}\textbf{Section 9.5: Basics of Cryptography}}\\*\hline
\textbf{9.5.1 Secret-key cryptography} & The same secret or closely related secrets encrypt and decrypt. It is efficient for bulk data, but communicating parties must establish and protect shared keys.\\\hline
\textbf{9.5.2 Public-key cryptography} & A public key can be distributed while a private key remains secret. The pair supports encryption or signature schemes but typically costs more than symmetric operations.\\\hline
\textbf{9.5.3 One-way functions} & A cryptographic hash maps arbitrary input to a fixed-size digest and should resist preimages and collisions. Password verification stores a salted derived value rather than plaintext.\\\hline
\textbf{9.5.4 Digital signatures} & A signer applies a private-key operation to a message digest; verifiers use the public key. Certificates bind keys to identities through signed assertions and trust chains.\\\hline
\textbf{9.5.5 Trusted platform modules} & A TPM protects keys and records measurements so software state can be attested. Hardware support strengthens key custody but does not make unmeasured software secure.\\\hline
\rowcolor{ice}\multicolumn{2}{l}{\hypertarget{sec-9-6}{}\textbf{Section 9.6: Authentication}}\\*\hline
\textbf{Knowledge factors} & Passwords are easy to deploy but vulnerable to guessing, reuse, observation, and compromised verifiers. Salting, slow derivation, rate limits, and secure reset paths reduce risk.\\\hline
\textbf{9.6.1 Physical objects} & Possession factors include cards, hardware tokens, and one-time-code devices. Loss, cloning, enrollment, and recovery remain part of the threat model.\\\hline
\textbf{9.6.2 Biometrics} & Biometrics compare noisy measurements rather than exact secrets. Thresholds trade false acceptance against false rejection, and compromised biometric traits are difficult to replace.\\\hline
\textbf{Multiple factors} & Combining independent factor types increases assurance only when an attacker cannot compromise them through the same path.\\\hline
\rowcolor{ice}\multicolumn{2}{l}{\hypertarget{sec-9-7}{}\textbf{Section 9.7: Exploiting Software}}\\*\hline
\textbf{Exploit chain} & An attacker turns a programming error into unintended reads, writes, control flow, or command execution. Privilege and reachable inputs determine impact.\\\hline
\textbf{9.7.1 Buffer overflow} & Unchecked writes can corrupt adjacent data, return addresses, or function pointers. Bounds checks, memory-safe languages, stack canaries, nonexecutable data, and address randomization provide complementary barriers.\\\hline
\textbf{9.7.2 Format strings} & Using attacker-controlled text as a formatting directive can disclose memory or write through special conversions. Programs must use constant formats and pass input only as data arguments.\\\hline
\textbf{9.7.3 Dangling pointers} & After memory is freed, stale references may access a reallocated object. An attacker can shape reuse so a later dereference reads or invokes attacker-controlled data.\\\hline
\textbf{9.7.4 Null dereferences} & A null reference normally crashes a process, but historical systems that allowed low-address mappings could turn some kernel dereferences into controlled memory access.\\\hline
\textbf{9.7.5 Integer overflow} & Arithmetic wraparound or signedness mistakes can make allocation or bounds calculations smaller than the data later copied. Validation must precede narrowing and size computation.\\\hline
\textbf{9.7.6 Command injection} & Concatenating untrusted input into a shell or interpreter command lets data become syntax. Structured APIs, allowlists, separation of arguments, and least privilege reduce exposure.\\\hline
\textbf{9.7.7 TOCTOU} & A resource can change between a security check and later use. Atomic operations, stable handles, directory-relative calls, or locking must bind validation to the object actually used.\\\hline
\rowcolor{ice}\multicolumn{2}{l}{\hypertarget{sec-9-8}{}\textbf{Section 9.8: Insider Attacks}}\\*\hline
\textbf{9.8.1 Logic bombs} & Malicious code waits for a condition or time before triggering. Reviews, separation of duties, change control, and monitoring make hidden triggers harder to introduce.\\\hline
\textbf{9.8.2 Back doors} & A back door bypasses normal authentication or authorization. It may be intentionally planted or left as a debugging shortcut.\\\hline
\textbf{9.8.3 Login spoofing} & A counterfeit login interface captures credentials and may then display an error or start the real session. Trusted attention paths help users reach genuine authentication code.\\\hline
\rowcolor{ice}\multicolumn{2}{l}{\hypertarget{sec-9-9}{}\textbf{Section 9.9: Malware}}\\*\hline
\textbf{9.9.1 Trojan horses} & A Trojan presents a useful or expected function while performing hidden malicious actions. It relies on a user or administrator granting execution.\\\hline
\textbf{9.9.2 Viruses} & A virus attaches to host content and replicates when that host executes or is processed. Infection strategies vary by executable, boot, document, and polymorphic behavior.\\\hline
\textbf{9.9.3 Worms} & A worm is a self-contained network-spreading program. Automated discovery and exploitation can create rapid propagation and large-scale resource consumption.\\\hline
\textbf{9.9.4 Spyware} & Spyware monitors activity or collects information without informed consent, including keystrokes, browsing, credentials, or user behavior.\\\hline
\textbf{9.9.5 Rootkits} & A rootkit hides malicious presence or maintains privileged control by modifying kernel, driver, library, boot, or firmware behavior.\\\hline
\rowcolor{ice}\multicolumn{2}{l}{\hypertarget{sec-9-10}{}\textbf{Section 9.10: Defenses}}\\*\hline
\textbf{Defense in depth} & No single control stops every path. Prevention, least privilege, isolation, detection, recovery, patching, and secure administration should fail independently.\\\hline
\textbf{9.10.1 Firewalls} & Packet filters and stateful or application gateways enforce network traffic policy at boundaries. They cannot fully protect trusted protocols, encrypted malicious content, or internal paths.\\\hline
\textbf{9.10.2 Antivirus} & Scanners use signatures, heuristics, behavior, emulation, and reputation. Malware can obfuscate or detect analysis, so scanning is one layer rather than proof of safety.\\\hline
\textbf{9.10.3 Code signing} & A verified signature establishes that code matches what a recognized signer approved. Trust still depends on key protection, signer policy, and the signed code's quality.\\\hline
\textbf{9.10.4 Jailing} & A jail restricts namespace, resources, or system calls available to a process. Containment limits consequences but shares some trusted kernel surface.\\\hline
\textbf{9.10.5 Model-based detection} & An intrusion detector compares observed events with expected behavior or attack patterns. Sensitivity trades missed attacks against false alarms.\\\hline
\textbf{9.10.6 Mobile-code encapsulation} & Sandboxing, interpretation, capability restriction, or proof checking confines downloaded code so it cannot freely access host resources.\\\hline
\textbf{9.10.7 Java security} & \textcolor{legacy}{\textbf{Historical or legacy context:}} The model combines bytecode verification, class loading, type safety, and permission checks to constrain untrusted applets or components.\\\hline
\rowcolor{ice}\multicolumn{2}{l}{\hypertarget{sec-9-11}{}\textbf{Section 9.11: Research on Security}}\\*\hline
\textbf{Research themes} & Research spans verified kernels, exploit mitigations, authentication, information flow, malware analysis, hardware roots of trust, anomaly detection, and usable security.\\\hline
\rowcolor{ice}\multicolumn{2}{l}{\hypertarget{sec-9-12}{}\textbf{Section 9.12: Summary}}\\*\hline
\textbf{Layered assurance} & Security depends on a small trustworthy enforcement base, explicit access policy, sound authentication and cryptography, robust software, containment, monitoring, and recovery.\\\hline
\end{longtable}
\normalsize

\section*{Chapter Synthesis}
\begin{studybox}{Chapter Summary}
Operating-system security protects confidentiality, integrity, and availability through trusted mediation, identity, authorization, isolation, and cryptographic support. Software defects, insiders, and malware exploit gaps in these controls. Defense in depth combines prevention, containment, detection, and recovery because every individual mechanism has limits.
\end{studybox}

\begin{studybox}{Key Relationships}
\begin{itemize}
  \item Authentication establishes an identity; authorization maps that identity or domain to allowed operations.
  \item ACLs and capabilities are column- and row-oriented representations of the same abstract access matrix.
  \item Cryptography protects data and identity assertions, while the operating system protects keys, processes, and access paths.
  \item Exploit mitigations reduce reliability of individual bugs; sandboxing and least privilege reduce post-exploit authority.
\end{itemize}
\end{studybox}

\begin{studybox}{Design Tradeoffs}
\begin{itemize}
  \item A smaller TCB is easier to analyze, but moving services out can increase communication and interface complexity.
  \item Strong authentication improves assurance but can reduce usability and complicate recovery.
  \item Detection tuned for sensitivity finds more attacks but produces more false positives.
  \item Revocable centralized ACLs are easy to audit per object; capabilities make delegation and least authority natural.
\end{itemize}
\end{studybox}

\begin{studybox}{Common Confusions}
\begin{itemize}
  \item Authentication answers who or what is acting; authorization answers what that actor may do.
  \item Protection enforces resource access rules; security also addresses malicious behavior, assurance, monitoring, and recovery.
  \item Encryption provides confidentiality, a hash provides a digest, and a digital signature provides origin and integrity evidence under a key trust model.
  \item A virus needs a host, a worm spreads as a standalone program, and a Trojan relies on deceptive execution.
\end{itemize}
\end{studybox}

\section*{Key Terms}
\small
\begin{longtable}{C N}
\rowcolor{navy}\color{white}\textbf{Term} & \color{white}\textbf{Concise definition}\\
\endfirsthead
\rowcolor{navy}\color{white}\textbf{Term} & \color{white}\textbf{Concise definition (continued)}\\
\endhead
\textbf{Confidentiality} & Prevention of unauthorized disclosure.\\\hline
\textbf{Integrity} & Prevention or detection of unauthorized modification.\\\hline
\textbf{Availability} & Ability of authorized users to obtain timely service.\\\hline
\textbf{TCB} & Components whose correct behavior is required to enforce security.\\\hline
\textbf{Reference monitor} & Tamper-resistant, always-invoked access mediation concept.\\\hline
\textbf{Protection domain} & Current collection of rights under which execution proceeds.\\\hline
\textbf{ACL} & Object-associated list of subjects and allowed operations.\\\hline
\textbf{Capability} & Unforgeable object reference carrying access rights.\\\hline
\textbf{Covert channel} & Unintended resource used to transmit information.\\\hline
\textbf{Cryptographic hash} & Fixed-size digest designed for preimage and collision resistance.\\\hline
\textbf{Digital signature} & Private-key evidence verifiable with a corresponding public key.\\\hline
\textbf{Salt} & Unique value mixed into password derivation to defeat shared precomputation.\\\hline
\textbf{TOCTOU} & Race between checking a resource and using it.\\\hline
\textbf{Rootkit} & Mechanism that hides or preserves privileged malicious control.\\\hline
\textbf{Sandbox} & Restricted execution environment limiting accessible resources.\\\hline
\textbf{Defense in depth} & Multiple independent controls across prevention, detection, and recovery.\\\hline
\end{longtable}
\normalsize

\enlargethispage{2\baselineskip}
\begin{studybox}{Active-Recall Review}
\small
\begin{enumerate}
  \item How do confidentiality, integrity, and availability differ?
  \item Why should the trusted computing base be small?
  \item How do ACLs and capabilities represent the access matrix?
  \item What makes a covert channel different from an ordinary permitted channel?
  \item How do symmetric encryption, public-key cryptography, hashing, and signatures serve different goals?
  \item What tradeoff does biometric threshold selection create?
  \item How can an integer overflow become a memory corruption vulnerability?
  \item Why is TOCTOU fundamentally an object-identity and atomicity problem?
  \item How do viruses, worms, Trojan horses, spyware, and rootkits differ?
  \item Why are code signing and malware scanning insufficient alone?
\end{enumerate}
\end{studybox}

\par\medskip
{\color{navy}\bfseries Next-Chapter Connections}\par\smallskip
Chapter 10 shows how process, memory, I/O, file-system, and security mechanisms are integrated in UNIX, Linux, and the historically described Android platform.
\normalsize
\end{document}
